% Shared typesetting and notation for the two associated documents.
\usepackage[T1]{fontenc}
\usepackage[utf8]{inputenc}
\usepackage[english]{babel}
\usepackage{lmodern,microtype}
\usepackage[a4paper,margin=24mm,headheight=14pt]{geometry}
\usepackage{amsmath,amssymb,amsthm,mathtools,mathrsfs}
\usepackage{booktabs,longtable,array,tabularx}
\usepackage{xcolor,enumitem,fancyhdr,needspace,etoolbox,placeins}
\usepackage{xr-hyper}
\usepackage[colorlinks=true,linkcolor=blue!40!black,citecolor=blue!40!black,urlcolor=blue!40!black]{hyperref}
\usepackage{bookmark,xurl}
\newcommand{\E}{\mathbb E}
\renewcommand{\P}{\mathbb P}
\newcommand{\N}{\mathbb N}
\newcommand{\Z}{\mathbb Z}
\newcommand{\F}{\mathbb F}
\newcommand{\Law}{\mathcal L}
\newcommand{\Pois}{\operatorname{Pois}}
\newcommand{\TV}{d_{\mathrm{TV}}}
\newcommand{\ind}{\mathbf 1}
\newcommand{\eps}{\varepsilon}
\DeclareMathOperator{\Var}{Var}
\DeclareMathOperator{\Cov}{Cov}
\DeclareMathOperator{\rank}{rank}
\DeclareMathOperator{\codim}{codim}
\renewcommand{\sf}{\operatorname{sf}}
\newtheorem{theorem}{Theorem}[section]
\newtheorem{proposition}[theorem]{Proposition}
\newtheorem{lemma}[theorem]{Lemma}
\newtheorem{corollary}[theorem]{Corollary}
\theoremstyle{definition}
\newtheorem{definition}[theorem]{Definition}
\theoremstyle{remark}
\newtheorem{remark}[theorem]{Remark}
\newtheorem{example}[theorem]{Example}
\numberwithin{equation}{section}
\setlength{\parindent}{1.2em}
\setlength{\parskip}{3pt}
\setlength{\emergencystretch}{2em}
\setlist{itemsep=3pt,topsep=5pt,leftmargin=20pt}
\renewcommand{\arraystretch}{1.12}
\pagestyle{fancy}
\fancyhf{}
\fancyhead[L]{\small Long multiplicative runs}
\fancyhead[R]{\small Brice Pouly}
\fancyfoot[C]{\thepage}
\renewcommand{\headrulewidth}{0.3pt}
\setcounter{tocdepth}{1}
\hypersetup{
 pdftitle={Sharp Poisson approximation and second-order correction for long multiplicative runs},
 pdfauthor={Brice Pouly},
 pdfsubject={Random multiplicative functions, Poisson approximation, hereditary stopping, arithmetic dependencies},
 pdfkeywords={Random multiplicative functions; long runs; optimal Poisson approximation; factorial cumulants; working draft 1.8, 7 November 2026},
 bookmarksdepth=2}
\allowdisplaybreaks[1]
\widowpenalty=10000
\clubpenalty=10000

% Presentation-only identifiers for the separate condensed preview.
\fancyfoot[L]{\scriptsize\color{black!60}Unified editorial preview}
\hypersetup{pdftitle={Long multiplicative runs: unified editorial preview},pdfkeywords={Unified editorial preview; not an adopted manuscript revision}}
